\begin{titlepage}
\thispagestyle{plain}
\centering
\vspace*{0.15in}
{\sffamily\bfseries\fontsize{17.28}{21}\selectfont \booktitle\par}
\vspace{5pt}
{\itshape \booksubtitle\par}
\vspace{10pt}
\bookauthor\par
\vspace{10pt}
Version: \bookversion\par
\vspace{24pt}
{\fboxsep=0pt\fboxrule=3pt\fbox{% Cover art: a Napkin-style collage (watercolor washes, spheres, hatched streaks, plus marks, a lettered sticker).
% Drawn with a fixed pseudo-random seed, so the picture is identical on every build. Replace or add motifs that
% belong to the book's topic (a robotics book might draw a robot arm, a cup, a camera lens and a speech bubble); keep the
% sticker text short enough to fit (check the rendered title page).
\begin{tikzpicture}[x=0.55in,y=0.55in]
  \pgfmathsetseed{20261103}
  \clip (0,0) rectangle (10,10);
  % paper and faint graph lines
  \fill[teal!8!white] (0,0) rectangle (10,10);
  \foreach \i in {1,...,19} {
    \draw[teal!28!white,line width=0.25pt] (\i*0.5,0) -- (\i*0.5,10);
    \draw[teal!28!white,line width=0.25pt] (0,\i*0.5) -- (10,\i*0.5);
  }
  % big washes
  \shade[inner color=violet!35!white,outer color=violet!70!blue!85,opacity=0.95] (1.9,7.7) circle (2.55);
  \shade[inner color=teal!40!white,outer color=teal!85!cyan!95,opacity=0.95] (1.2,1.4) circle (2.95);
  \shade[inner color=magenta!15!white,outer color=magenta!45!white,opacity=0.95] (8.4,1.6) circle (2.7);
  \shade[inner color=yellow!40!white,outer color=yellow!85!orange,opacity=0.9] (3.3,4.6) circle (1.2);
  % spheres
  \shade[ball color=orange!90!yellow] (6.75,7.0) circle (1.25);
  \shade[ball color=violet!70!blue] (8.0,8.35) circle (1.35);
  \shade[ball color=red!75!black] (0.9,3.35) circle (0.5);
  % hatched streaks
  \foreach \i in {0,...,52} {
    \pgfmathsetmacro{\hh}{1.9+1.6*rnd}
    \draw[magenta!60!purple,opacity=0.45,line width=0.9pt] (6.4+\i*0.07,0) -- ++(0,\hh);
  }
  \foreach \i in {0,...,22} {
    \pgfmathsetmacro{\hh}{0.9+0.9*rnd}
    \draw[blue!65!cyan,opacity=0.3,line width=0.9pt] (3.0+\i*0.06,9.9) -- ++(0.9*\hh,-0.7*\hh);
  }
  % small watercolor strokes
  \foreach \i in {1,...,160} {
    \pgfmathsetmacro{\sx}{10*rnd}
    \pgfmathsetmacro{\sy}{10*rnd}
    \pgfmathsetmacro{\sa}{360*rnd}
    \draw[white,opacity=0.10,line width=2.2pt] (\sx,\sy) -- ++(\sa:0.7);
  }
  % plus marks, as on the Napkin
  \foreach \i in {0,...,5} {
    \foreach \j in {0,...,2} {
      \draw[white,line width=0.5pt,opacity=0.8] (0.45+\i*0.55,1.55+\j*0.45) -- ++(0.14,0) (0.52+\i*0.55,1.48+\j*0.45) -- ++(0,0.14);
    }
  }
  % a band of decaying waves (replace with a motif of the topic)
  \foreach \k/\sc in {0/crimson,1/blue!75!black,2/orange!90!yellow,3/green!55!black,4/violet,5/cyan!70!blue} {
    \pgfmathsetmacro{\ph}{\k*37}
    \draw[\sc,line width=1.1pt,opacity=0.9] plot[domain=0:10,samples=140,smooth]
      ({\x},{2.6+0.12*\k+1.3*exp(-0.25*\x)*sin(95*\x+\ph)});
  }
  % sticker with the title
  \begin{scope}[rotate around={-4:(5,5)}]
    \fill[white,opacity=0.94,rounded corners=14pt] (1.4,4.1) rectangle (8.6,6.3);
    \draw[black,line width=1.4pt,rounded corners=14pt] (1.4,4.1) rectangle (8.6,6.3);
    \foreach \a in {0,45,...,315} {
      \node[font=\sffamily\bfseries\fontsize{30}{32}\selectfont,text=black] at (5+0.045*cos \a,5.35+0.045*sin \a) {\bookshorttitle};
    }
    \node[font=\sffamily\bfseries\fontsize{30}{32}\selectfont,text=white] at (5,5.35) {\bookshorttitle};
    \node[font=\itshape\fontsize{10.5}{12}\selectfont,text=black!70] at (5,4.45) {a study book};
  \end{scope}
\end{tikzpicture}
}\par}
\end{titlepage}

% verso: epigraph, colophon
\newpage
\thispagestyle{empty}
\vspace*{0.14\textheight}
\hfill\begin{minipage}{0.66\textwidth}
\raggedleft\itshape
% A short quotation that fits the topic. Check the wording against a primary source (or a source that quotes
% it verbatim) and cite it; say in provenance/README.md how it was checked.
``\fillin{epigraph}''\par
\medskip
\upshape\raggedleft---\,\fillin{its author, with a citation}\par
\end{minipage}
\vfill
\noindent\rule{\textwidth}{0.6pt}\par
\smallskip
\begingroup\small\parindent=0pt\parskip=4pt
Written for readers who \fillin{what the reader already knows, and what this book does not assume}.

Compiled \bookdate, with Claude Code. \fillin{what was verified, against what, and when: for example, the titles,
authors, dates and categories of the arXiv papers were checked against the arXiv API on <date>}. The technical
descriptions are the compiler's synthesis of the cited sources, and independent fact-checks of the text are kept
with the sources (see the Preface).

% Keep this paragraph only if papers are appended to the combined edition.
The papers reproduced in Part~\ref{part:papers} are the authors' works, as distributed by arXiv under
each paper's own terms. This volume is a personal study copy.

Version: \bookversion.\par
\endgroup
\vspace*{0.3in}
\newpage
